\documentclass[11pt]{article}
\usepackage[T1]{fontenc}
\usepackage[utf8]{inputenc}
\usepackage{lmodern}
\usepackage[margin=1in]{geometry}
\usepackage{amsmath,amssymb,amsthm,mathtools}
\usepackage{microtype}
\usepackage[hidelinks]{hyperref}
\setlength{\emergencystretch}{2em}
\newtheorem{theorem}{Theorem}[section]
\title{Sharp Four-State Hierarchy Prices\protect\\for Arbitrary Positive Weights}
\author{Phase 9 research verification copy}
\date{7 October 2026}
\begin{document}
\maketitle
\noindent\textbf{Status.} Internal mathematical review, not external peer
review or a claim of historical priority. The included theorem concerns
exactly four states and output budgets two and three.

\section{Previously established lower bound}
\begin{theorem}[Phase 8 endpoint lower bound]
\label{thm:ext-weightdegenerate}
There are absolute $c_*>0$ and $r_*>0$ such that, for every integer
$r\ge r_*$, four-state sources with positive weights and nonnegative
polynomial curvature of degree at most $r$ have both deterministic and
stochastic two-/three-output hierarchy prices at least
$c_*r^2/(\log r)^2$. The posteriors may lie in $[1/4,3/4]$, and the
generators may induce smooth strictly proper losses in $[0,1/2]$.
\end{theorem}
\noindent This is an imported statement only. Its complete proof and
normalization are in the immutable Phase 8 manuscript file
\texttt{unequal\_weights\_endpoint.tex}; its full independent review is
\texttt{unequal\_weights\_independent\_review.md}. All new upper-bound
calculations are given below.

% Phase 9 research candidate; include-ready for curvature_hierarchy.tex.
% All Phase 8 source files are immutable and were only read.
% Independently checked algebraically by endpoint_obstruction and
% phase9_adversarial on 7 October 2026. Internal review only.
\section{The sharp degree order for four arbitrary positive weights}
\label{sec:weighted-sharp}

The logarithmic gap in the four-state arbitrary-weight problem can be
closed on the upper-bound side. The proof uses the original source weights
in every Jensen kernel. Choosing the appropriate one of the two adjacent
triples is essential: a corresponding assertion for an arbitrary weighted
triple would be false.

\begin{theorem}[Weighted four-state degree bound]
\label{thm:weighted-sharp}
Let $g=\varphi''$ be a nonzero polynomial of degree at most $r$, nonnegative
on $[0,1]$. Let four source states have arbitrary positive probabilities
and posterior locations in $[0,1]$, and suppose that both flat distortions
$\operatorname{OPT}_2$ and $\operatorname{OPT}_3$ are positive. Write
$R=\rho_{\varphi,\mathrm{det}}$ for the optimal deterministic hierarchy
price at output budgets two and three.
If $r\ge1$ and $R\ge128$, then
\begin{equation}
 R^2\le16r^2\exp\left(\frac{20r}{\sqrt R}\right).
 \label{eq:weighted-implicit}
\end{equation}
In particular, if $R\ge256$ then
\begin{equation}
 R(\log R)^2\le1936r^2.
 \label{eq:weighted-Rlog}
\end{equation}
For every integer $r\ge2$ this gives the explicit uniform bound
\begin{equation}
 \rho_{\varphi,\mathrm{stoch}}
 \le\rho_{\varphi,\mathrm{det}}
 \le1936\frac{r^2}{(\log r)^2}.
 \label{eq:weighted-sharp-upper}
\end{equation}
Consequently, in the four-state degree envelopes permitting all positive
weights and all posterior locations, both prices have sharp order
\[
 \widehat C_r^{\mathrm{det}}
 =\Theta\!\left(\frac{r^2}{(\log r)^2}\right),\qquad
 \widehat C_r^{\mathrm{stoch}}
 =\Theta\!\left(\frac{r^2}{(\log r)^2}\right).
\]
The matching lower bounds are the endpoint construction of
Theorem~\ref{thm:ext-weightdegenerate}.
\end{theorem}

\begin{proof}
We first record the two elementary polynomial estimates used below. If
$f\ge0$ is a polynomial of degree at most $r\ge1$ on an interval $E$ of
length $L>0$, and $M_E=\max_E f$, then the rescaled Markov inequality gives
\begin{equation}
 \int_E f(t)\,\mathrm dt\ge\frac{L M_E}{8r^2}.
 \label{eq:weighted-markov-mass}
\end{equation}
Indeed, $\|f'\|_\infty\le2r^2M_E/L$. From a maximizer there is an
available one-sided segment of length $L/(4r^2)$, and on this segment
$f\ge M_E/2$.

If $t_0$ lies to the right of $E$, at distance $d$ from its right endpoint,
the classical Chebyshev extremal inequality gives
\begin{equation}
 |f(t_0)|\le M_E T_r(1+2d/L)
 \le M_E\exp\left(2r\sqrt{d/L}\right).
 \label{eq:weighted-cheb-extrapolation}
\end{equation}
Here the second estimate follows from
$T_r(\cosh s)=\cosh(rs)$ and
$\operatorname{arcosh}(1+x)\le\sqrt{2x}$. To see the first estimate
directly, rescale $E$ to $[-1,1]$ and divide by $M_E$. If a degree-at-most-$r$
polynomial $F$ bounded by one there had
$|F(z)|>T_r(z)$ at $z>1$, change its sign if necessary and set
$\lambda=T_r(z)/F(z)\in(0,1)$. At the $r+1$ alternating extrema of $T_r$,
the polynomial $T_r-\lambda F$ has the same strict alternating signs as
$T_r$, giving $r$ roots in $[-1,1]$, as well as its root at $z$.
This contradicts its degree. These estimates are classical polynomial
tools, not new approximation results.

\emph{The sole optimizer conflict is still the middle pair.}
Positive $\operatorname{OPT}_3$ implies that the four posteriors are
distinct; label them $A<B<C<D$. For a cell or pair, retain its original
source weights without conditional renormalization. Put
\[
 p_L=D(AB),\quad q=D(BC),\quad p_R=D(CD),\quad
 Q=p_L+p_R,\quad T_L=D(ABC),\quad T_R=D(BCD).
\]
Contiguous flat optimizers exist for arbitrary positive weights: for fixed
centers $a<b$, the difference $B_\varphi(x,a)-B_\varphi(x,b)$ is affine
in $x$ with positive slope, so nearest-center assignments are contiguous,
and updating to weighted centroids cannot increase distortion. Strict
convexity and distinct posteriors ensure that the optimum fills its budget.

For every weighted triple, its kernel dominates the two constituent
adjacent-pair kernels on their respective disjoint interiors. One way to
see the domination is to merge a pair first and then merge its weighted
centroid with the remaining state; the additional Jensen gap is
nonnegative, including for each hinge generator. Thus
\[
 T_L\ge p_L+q,\qquad T_R\ge q+p_R.
\]
If $AB$ is a cheapest adjacent pair, the only incompatible contiguous
coarse candidate is $A\mid BCD$; its cost obeys
$T_R\ge q+p_R\ge p_L+p_R=Q$, so a compatible flat-optimal coarse
partition can be chosen. The case of $CD$ is its reflection. If $BC$ is
cheapest and either triple split is flat optimal, it too has compatible
flat optima. Hence either $R=1$, or we may work with the sole conflict
\[
 \operatorname{OPT}_3=q,\qquad \operatorname{OPT}_2=Q,
\]
where the optimal fine pair is $BC$ and the optimal coarse split is
$AB\mid CD$.

The hierarchy retaining that coarse split and its cheaper outer fine pair,
and the two hierarchies retaining the fine pair $BC$ and an adjacent triple,
give the admissible-candidate inequalities
\begin{equation}
 R\le\frac{\min(p_L,p_R)}q,\qquad
 R\le\frac{T_L}Q,\qquad R\le\frac{T_R}Q.
 \label{eq:weighted-candidates}
\end{equation}
They use feasible deterministic hierarchies only. They do not require an
exhaustion of optimal noncontiguous hierarchies.

\emph{Choose the triple by the middle masses.}
Reflect the ordered configuration if necessary so that $p_B\le p_C$.
Use the triple $BCD$, and affinely normalize its own hull to $[0,1]$:
\[
 B=0,\quad C=c\in(0,1),\quad D=1,\qquad
 (p_B,p_C,p_D)=(u,v,w),\quad 0<u\le v,\quad w>0.
\]
The fourth point is not used in the polynomial estimates. If the original
triple hull has length $L_0$, the transformed curvature is
$L_0^2g(B+L_0t)$, preserving the curvature degree, nonnegativity, and all
the costs under discussion. Denote it again by $g$.

Let $K$ be the triple kernel, let $k_q$ and $k_p$ be its two adjacent-pair
kernels, and let $H$ be the kernel of its two outer states with their
original weights:
\begin{align*}
 K(t)&=\min\{ut+v(t-c)_+,\ w(1-t)+v(c-t)_+\},\\
 k_q(t)&=\min\{ut,v(c-t)\}\,\mathbf1_{[0,c]}(t),\\
 k_p(t)&=\min\{v(t-c),w(1-t)\}\,\mathbf1_{[c,1]}(t),\\
 H(t)&=\min\{ut,w(1-t)\}.
\end{align*}
In particular $T:=T_R=\int_0^1Kg$, $q=\int_0^1k_qg$, and
$p_R=\int_0^1k_pg$. Write $S=p_R+q$.
By \eqref{eq:weighted-candidates},
\begin{equation}
 q\le\frac{Q}{2R}\le\frac{T}{2R^2},\qquad
 S\le Q+q\le\frac{T}{R}+\frac{T}{2R^2}\le\frac{2T}{R}.
 \label{eq:weighted-two-small-costs}
\end{equation}

Two pointwise facts make these estimates uniform in the weights. First,
\begin{equation}
 0\le K-H\le k_q+k_p.
 \label{eq:weighted-H-comparison}
\end{equation}
For $t<c$, the difference is
$\min\{ut,w(1-t)+v(c-t)\}-\min\{ut,w(1-t)\}$,
which lies between zero and $\min\{ut,v(c-t)\}=k_q(t)$.
For $t>c$, interchange the two sides to obtain the bound by $k_p(t)$.
Second, for $t<c$, $K(t)\le ut$ gives
\[
 \frac{k_q(t)}{K(t)}
 \ge\min\left\{1,\frac{v(c-t)}{ut}\right\}
 \ge\min\left\{1,\frac{c-t}{t}\right\}.
\]
For $t>c$, put $\Delta=t-c$. The elementary inequality
$\min\{B,C\}/\min\{A+B,C\}\ge B/(A+B)$, for positive
$A,B,C$, gives
\[
 \frac{k_p(t)}{K(t)}
 \ge\frac{v\Delta}{ut+v\Delta}
 \ge\frac{\Delta}{c+2\Delta}.
\]
Endpoint values follow by continuity; ratios need only be used inside
$(0,1)$, where $K>0$.

\emph{A large peak and a nearby interval with small curvature cost.}
Assume $R\ge128$ and put
\[
 \varepsilon=\frac{16}{R}\le\frac18,
 \qquad I=\{t\in[0,1]:|t-c|\le\varepsilon c\}.
\]
The preceding pointwise inequalities imply, outside $I$,
\[
 k_q+k_p\ge\frac{\varepsilon}{1+2\varepsilon}K.
\]
Consequently
\[
 \int_{[0,1]\setminus I}Kg
 \le\frac{1+2\varepsilon}{\varepsilon}S
 \le\frac{5T}{32}.
\]
By \eqref{eq:weighted-H-comparison} and
\eqref{eq:weighted-two-small-costs},
\[
 \int_I Hg\ge T-\frac{5T}{32}-S
 \ge\frac{53T}{64}\ge\frac{3T}{4}.
\]
On $I$, $H(t)\le ut\le u(1+\varepsilon)c$, and $|I|\le2\varepsilon c$.
Thus there exists $t_0\in I$ such that
\begin{equation}
 g(t_0)\ge\frac{T}{4\varepsilon u c^2}.
 \label{eq:weighted-peak}
\end{equation}
This estimate uses $H\le ut$, rather than dividing by $H(c)$, which could
be arbitrarily small as weights or gaps degenerate.

Consider the interval
\[
 E=[c/4,c(1-2\varepsilon)].
\]
Its length is at least $c/2$. The point $t_0$ lies to its right at distance
at most $3\varepsilon c$. On all of $E$, both terms in
$k_q(t)=\min\{ut,v(c-t)\}$ are at least $2u\varepsilon c$:
indeed $t\ge c/4\ge2\varepsilon c$, $c-t\ge2\varepsilon c$, and
$v\ge u$. With $M_E=\max_Eg$, the extrapolation estimate
\eqref{eq:weighted-cheb-extrapolation} and \eqref{eq:weighted-peak} give
\[
 M_E\ge g(t_0)\exp(-\sqrt{24}\,r\sqrt\varepsilon)
 \ge\frac{T}{4\varepsilon u c^2}
       \exp\left(-\frac{20r}{\sqrt R}\right),
\]
where $4\sqrt{24}<20$. The mass estimate
\eqref{eq:weighted-markov-mass} now yields
\begin{align*}
 q&\ge2u\varepsilon c\int_Eg(t)\,\mathrm dt\\
  &\ge2u\varepsilon c\frac{cM_E}{16r^2}
   \ge\frac{T}{32r^2}
       \exp\left(-\frac{20r}{\sqrt R}\right).
\end{align*}
Combining this with $q\le T/(2R^2)$ proves
\eqref{eq:weighted-implicit}.

\emph{Extract the degree order.}
Set $z=r/\sqrt R>0$. Equation~\eqref{eq:weighted-implicit} implies
\[
 \log(R/16)\le20z+2\log z\le22z.
\]
For $R\ge256$, $\log(R/16)\ge\tfrac12\log R$, so
$\sqrt R\log R\le44r$, proving \eqref{eq:weighted-Rlog}.
For $r\ge2$, if $R\ge256$ and $R>r$, substitute
$\log R\ge\log r$ to obtain \eqref{eq:weighted-sharp-upper}.
If $R\le r$, use $(\log r)^2\le r$; if $R<256$, use
$(\log r)^2\le r^2$ and $256<1936$.
Every deterministic hierarchy is stochastic admissible, with the same flat
denominators, so the stochastic upper bound follows. The already proved
endpoint lower construction gives the two matching lower orders.
\end{proof}

\paragraph{Quantifiers, normalization, and boundaries.}
The constant is independent of all positive source probabilities and all
four distinct posterior locations. The proof does not impose a lower bound
on a posterior gap or a source mass. All cell kernels use the original
source weights; the selected triple is never conditionally renormalized.
Only its posterior coordinate is rescaled, with the corresponding squared
length factor in the curvature.

The matching lower witnesses can be placed in $[1/4,3/4]$ and represented
by strictly proper losses in $[0,1/2]$, exactly as in
Theorem~\ref{thm:ext-weightdegenerate}. For its notation
$d=(4\log n/n)^2$, $b=d^2$, and
$M_n=\max_{[-1/2,3/2]}g_n$, the bounded-loss affine lift multiplies every
distortion by $1/(4M_n)$. The explicit feasible flat bounds therefore are
\[
 0<\widetilde{\operatorname{OPT}}_3
 \le\frac{d^4}{4M_n},\qquad
 0<\widetilde{\operatorname{OPT}}_2
 \le\frac{2d^3}{M_n}.
\]
Since $M_n\ge\int_0^1g_n=2+d^2$, both upper bounds tend to zero.
The sharp hierarchy-price statement is multiplicative in excess risk;
it does not supply nonvanishing absolute excess costs after bounded-loss
normalization. No loss normalization is needed for the upper bound.

This theorem closes the gap only for four source states and budgets two
and three. On four states the all-capacity problem has the same optimum,
because the one-output and four-output levels are respectively the constant
representation and the identity. The proof does not improve the existing
$O(r^2)$ theorem for arbitrary finite source counts and two prescribed
capacities, nor the $O(r^2\log(r+2))$ theorem for one hierarchy covering all
capacities on arbitrary finite sources. It makes no claim of optimal
leading constants or historical priority.

\end{document}
